% ======================================================================
% CHAPTER: Abstract + Keywords (not numbered, not in ToC)
% ----------------------------------------------------------------------
% NOTE:
% \begin{chapterabstractpage} 
% This adds an abstract. 
% Also I set up the chapter to be an experimental chapter with the
% appropriate sections, if this is not needed please remove.
% ======================================================================

\begin{chapterabstractpage}
  % --- Write a concise abstract: background, objective, methods, results, conclusions.
  Briefly state background, objective, methods, main results, and conclusions.

  \bigskip
  % --- Keywords (plain text; adjust to your conventions if needed)
  \textbf{Keywords:} LaTeX; Overleaf; doctoral templates; margin alignment; header standardization; citation management; iterative despair; typographic integrity
\end{chapterabstractpage}


% ======================================================================
% MAIN BODY
% ======================================================================

% ----------------------------------------------------------------------
% 1. INTRODUCTION
%    - State the problem, why it matters, prior work, and research questions.
% ----------------------------------------------------------------------
\section{Introduction}
State the problem, why it matters, what is known, and your research question(s).


% ----------------------------------------------------------------------
% 2. METHODS
%    - Structure: Design → Participants/Data → Measures/Variables → Analysis
%    - Keep each subsection brief and to the point.
% ----------------------------------------------------------------------
\section{Methods}

% -- 2.1 Design
%    Include preregistration, ethics, deviations from protocol.
\subsection{Design}
Study design, preregistration, ethics, deviations.

% -- 2.2 Participants / Data
%    Sampling frame, inclusion/exclusion criteria, data sources.
\subsection{Participants / Data}
Sample, inclusion/exclusion, data sources.

% -- 2.3 Measures / Variables
%    Operationalizations, instruments, reliability/validity.
\subsection{Measures / Variables}
Operationalizations, instruments, reliability.

% -- 2.4 Analysis
%    Statistical/qualitative methods, software, thresholds.
\subsection{Analysis}
Statistical/qualitative methods, software, thresholds.


% ----------------------------------------------------------------------
% 3. RESULTS
%    - Present findings clearly.
%    - Refer to tables/figures in text (e.g., “see Table~\ref{tab:template_emotions}”).
% ----------------------------------------------------------------------
\section{Results}
Report findings clearly; use figures/tables as needed.

% ----- TABLE: “Emotional stages … Overleaf”
% Layout notes:
% - Uses tabularx to auto-wrap the Description column (type X).
% - Uses booktabs for clean rules (\toprule, \midrule, \bottomrule).
% - \arraystretch adjusted locally for row height; safe for layout.

\begin{table}[ht]
  \centering
  \renewcommand{\arraystretch}{1.1}
  \caption{Emotional stages of making a PhD thesis template in Overleaf.}
  \label{tab:template_emotions_Ch5}
  \begin{tabularx}{\textwidth}{@{}lXr@{}}
    \toprule
    \textbf{Phase} & \textbf{Description} & \textbf{Duration (hrs)} \\ \midrule
    Initial spark     & “This will be easy.”                      & 0.3 \\
    Mild confusion    & “Why no ToC entry?”                       & 0.8 \\
    Frustration spike & Missing \texttt{\textbackslash end\{document\}}? & 2.5 \\
    Existential doubt & “Word wasn’t that bad.”                   & 1.5 \\
    Brief euphoria    & “It finally compiles!”                    & 0.2 \\
    Design obsession  & “Adjust margins again.”                   & 4.0 \\
    False confidence  & “Looks perfect now.”                      & 0.0 \\
    Regression loop   & “Why broken today again?”                 & $\infty$ \\
    \bottomrule
  \end{tabularx}
\end{table}

% ----- FIGURE: Example figure
% Notes:
% - Place your file at the given path or adjust the path accordingly.
% - Width uses \linewidth fraction to scale relative to the text block.
\begin{figure}[ht]
  \centering
  \includegraphics[width=.8\linewidth]{figures/Chapter5/Figure.jpg}
  \caption{Example caption.}
  \label{fig:paper4:fig1}
\end{figure}


% ----------------------------------------------------------------------
% 4. DISCUSSION
%    - Interpret findings, discuss strengths/limits, and outline implications.
% ----------------------------------------------------------------------
\section{Discussion}

% -- 4.1 Principal Findings
%    High-level interpretation: what did you learn and why does it matter?
\subsection{Principal Findings}
Summarize key results and interpret them.

% -- 4.2 Strengths and Limitations
%    Be specific about internal/external validity; acknowledge constraints.
\subsection{Strengths and Limitations}
What’s robust, what’s not.

% -- 4.3 Implications and Future Work
%    Practical/theoretical takeaways + concrete next steps.
\subsection{Implications and Future Work}
Practical/theoretical implications; next steps.


% ======================================================================
% OPTIONAL PER-CHAPTER EXTRAS
% (Leave commented unless needed; uncommenting won’t change layout much,
%  but keep any additional commands/styles in the preamble.)
% ======================================================================

% -- Acknowledgements specific to this chapter (unnumbered)
% \section*{Acknowledgements}
% Thanks to ...

% -- Per-chapter references (biblatex)
% To have references at the end of THIS CHAPTER instead of global:
% 1) Wrap the entire chapter file in \begin{refsection} ... \end{refsection}
%    from main.tex, and
% 2) Uncomment the print line below.
% \printbibliography[heading=subbibliography,title={References for this chapter}]